\section{Discussion}
\label{sec:discussion}

The central payoff of this draft is conceptual clarity with enough concrete
mathematics behind it to be checkable. The paper does not merely rename
familiar notions. It shows that once perturbation type, evaluation point, and
deterministic tie-breaking are all fixed explicitly, one gets both a clean
uniform perturbation calculus and a sharp pointwise separation phenomenon for
nearest neighbors.

At the current stage, the strongest established claims are the following.
Replace-one admits a delete-plus-add upper bound at the uniform loss level
(\Cref{prop:replace-delete-add}). Deterministic $k$-NN prediction changes are
exactly margin-crossing events inside the ordered top-$k$ neighborhood
(\Cref{prop:margin-crossing}). And deterministic $1$-NN already exhibits a
finite metric separation between pointwise LOO stability and pointwise
replace-one stability (\Cref{prop:two-point-separation}).

The weaker parts of the story are equally important to state honestly. The
minimality of the current $1$-NN witness beyond the explicit construction is
supported by computational certificates, not by a full impossibility proof. The
odd-$k$ extension is presently a conjectural lifting pattern backed by bounded
search evidence for $k=3,5,7$. Those are meaningful research assets, but they
should not be promoted beyond their evidentiary level.

The draft also has several limitations that shape the next phase.
\begin{enumerate}[leftmargin=*,label=(\arabic*)]
    \item The current theorem-level separation is for deterministic finite
    $1$-NN; the general odd-$k$ story still lacks a formal proof.
    \item The ambient spaces are finite metric spaces, often realized as graph
    shortest-path metrics. This gives crisp witnesses but does not yet address
    broader geometric classes.
    \item The paper studies binary 0--1 loss. Real-valued scores, multiclass
    labels, and regression analogues remain open.
    \item The discussion is worst-case and fixed-sample. It is intentionally not
    a statistical rate paper.
    \item The computational certificates are reproducible and well documented,
    but they are still certificates rather than formal proofs.
\end{enumerate}

These limitations point directly to the most natural follow-up problems.
\begin{enumerate}[leftmargin=*,label=(\arabic*)]
    \item Prove or refute \Cref{conj:odd-k-lifting}.
    \item Strengthen the current certificate pipeline into a formal minimality
    argument or an automated proof assistant artifact.
    \item Characterize which metric-space constraints eliminate the separation.
    \item Extend the perturbation calculus to other local rules, such as
    weighted neighbors or radius-based classifiers.
    \item Connect the finite worst-case indicators here to distributional
    stability notions without collapsing the distinctions this paper isolates.
\end{enumerate}

If there is a broader moral, it is this: for local learning rules, the geometry
of the neighborhood and the semantics of the perturbation must be specified
together. Once that is done, some relationships become cleaner than expected,
and some purportedly similar notions separate much sooner than intuition
suggests.
